\begin{convention}[Index suppression and aggregation from the root variable]
\label{conv:notation}

Every variable in the model is defined with a fixed ordered tuple of
\textbf{root indices} $(a_1, a_2, \dots, a_k)$, each ranging over a
finite set $S_1, S_2, \dots, S_k$ respectively.
The full collection of root instances
\[
  \bigl\{\,\phi_{a_1 \cdots a_k} : a_\ell \in S_\ell,\ \ell = 1,\dots,k\bigr\}
\]
can be thought of as a table: each row is one tuple
$(a_1,\dots,a_k)\in S_1\times\cdots\times S_k$, each column
corresponds to one index, and the entry is the value of $\phi$ at
that tuple.
We call this the \textbf{root form} of $\phi$.

Writing a variable with a \emph{strict subset} of its root indices
always denotes a \textbf{derived form}: the missing indices have been
suppressed and their values aggregated according to the rules below.
Because the written indices identify exactly which columns have been
retained, the aggregate is unambiguous and no additional decoration
is required.

Figure~\ref{fig:index-lattice} illustrates the full derivation lattice.

% ------------------------------------------------------------------
\medskip
\noindent\textbf{Rule~1 (Sum aggregation).
  Applies to all variable types: continuous, integer, and binary.}

Suppressing one or more indices from $\phi$ corresponds to the
following operation on the root table: remove the columns for the
suppressed indices, then replace every group of rows that agree on
all remaining columns with a single row whose value is their sum.
The result is a new table indexed by the retained indices alone.

Formally, suppressing index $a_\ell$ yields the variable
$\phi_{a_1 \cdots a_{\ell-1}\, a_{\ell+1} \cdots a_k}$, defined by
%
\begin{equation*}
  \phi_{a_1 \cdots a_{\ell-1}\, a_{\ell+1} \cdots a_k}
  \;\coloneqq\;
  \sum_{a_\ell \in S_\ell}
  \phi_{a_1 \cdots a_k}.
\end{equation*}
%
Suppressing several indices simultaneously applies the same
operation once per suppressed column; the result is the same
regardless of the order in which columns are removed, and the
retained indices are written in their original relative order.

For continuous and integer variables this is the natural
interpretation.
For binary variables the result is an \emph{integer count} of how
many root-level binaries in the group equal~1; it is not itself
binary.
No new variable is introduced by Rule~1, and no constraint is
required.

% ------------------------------------------------------------------
\medskip
\noindent\textbf{Rule~2 (Dot notation for coincident index pairs).}

When two root indices range over the \emph{same} set $S_\ell = S_p$,
the columns they occupy in the root table are drawn from an identical
domain.
In that case, suppressing one while retaining the other is
ambiguous by position alone, since the surviving index name does not
reveal which of the two identical-domain columns was removed.
A dot~$(\,\cdot\,)$ placed in the position of the suppressed index
resolves this ambiguity explicitly:
%
\begin{equation*}
  \phi_{a_1 \cdots \cdot \cdots a_p \cdots a_k}
  \;\coloneqq\;
  \sum_{a_\ell \in S_\ell} \phi_{a_1 \cdots a_\ell \cdots a_p \cdots a_k},
\end{equation*}
%
where the dot occupies position $\ell$ and $a_p$ is retained.
The symmetric case, retaining position $\ell$ and suppressing
position $p$, is written analogously with the dot in position $p$.
Outside of such coincident pairs the written indices already
identify the retained columns unambiguously, so the dot is never
needed elsewhere.
Rule~1 applies as normal: the result is a sum, and for binary
variables it is an integer count.

% ------------------------------------------------------------------
\medskip
\noindent\textbf{Rule~3 (OR-aggregated binaries).
  A new binary variable derived from the sum aggregate.}

For binary variables, the sum aggregate produced by Rule~1 is an
integer count and is not itself binary.
When a binary indicator is needed at a coarser granularity---one
that signals whether \emph{at least one} root-level binary in the
group equals~1---a genuinely new binary decision variable is
introduced.
It is denoted by placing an overline over the variable symbol
($\overline{\phi}$) and carries the same retained indices as the
corresponding sum aggregate.

This variable is not a notational shorthand: it must be declared in
the solver and its relationship to the sum aggregate enforced by an
explicit big-$M$ pair.
Suppressing index $a_\ell$ from a binary variable $\phi$ yields the
sum aggregate
$\phi_{a_1 \cdots a_{\ell-1}\, a_{\ell+1} \cdots a_k}$
(an integer count, by Rule~1) and the new binary
$\overline{\phi}_{a_1 \cdots a_{\ell-1}\, a_{\ell+1} \cdots a_k}$,
linked by:
%
\begin{align}
  \overline{\phi}_{a_1 \cdots a_{\ell-1}\, a_{\ell+1} \cdots a_k}
    &\leq
    \phi_{a_1 \cdots a_{\ell-1}\, a_{\ell+1} \cdots a_k}
  \label{eq:or-lb} \\
  \phi_{a_1 \cdots a_{\ell-1}\, a_{\ell+1} \cdots a_k}
    &\leq
    M \cdot \overline{\phi}_{a_1 \cdots a_{\ell-1}\, a_{\ell+1} \cdots a_k}
  \label{eq:or-ub}
\end{align}
%
Constraint~\eqref{eq:or-lb} forces $\overline{\phi}$ to~0 when the
count is~0 (no root binary in the group is active).
Constraint~\eqref{eq:or-ub} forces $\overline{\phi}$ to~1 when the
count is positive (at least one root binary is active).
Together they enforce
$\overline{\phi} = 1 \iff \phi_{a_1 \cdots a_{\ell-1}\, a_{\ell+1} \cdots a_k} \geq 1$.

The same construction applies when several indices are suppressed
simultaneously: $\overline{\phi}$ is then indexed by all retained
columns, and the big-$M$ pair is stated in terms of the corresponding
multi-index sum aggregate.

% ------------------------------------------------------------------
\medskip
\noindent\textbf{Rule~4 (Exactly-one variant).
  An OR aggregate with an additional global restriction.}

In some cases the OR aggregates formed by Rule~3, viewed as a table
indexed by the retained indices, must satisfy the stronger condition
that \emph{exactly one} entry in each group of rows sharing a fixed
combination of all retained indices except $a_p$ equals~1.
Rather than introducing a further variable, the same OR aggregate
$\overline{\phi}$ from Rule~3 is reused and marked with a hat
($\hat{\phi}$) to signal that the following equality is imposed on
top of constraints~\eqref{eq:or-lb}--\eqref{eq:or-ub}:
%
\begin{align}
  \sum_{a_p} \hat{\phi}_{a_1 \cdots a_{\ell-1}\, a_{\ell+1} \cdots a_k}
    &= 1.
  \label{eq:exactly-one}
\end{align}
%
That is, $\hat{\phi}$ and $\overline{\phi}$ are the \emph{same
variable}---the OR aggregate of $\phi$ over $a_\ell$, constructed
identically via Rules~1 and~3.
The hat is purely a diacritic: it warns the reader that
constraint~\eqref{eq:exactly-one} further restricts the feasible
pattern of values across $a_p$.
Whereas plain OR aggregation allows any subset of rows to be
simultaneously active, the hat enforces that the sum of those OR
aggregates over all admissible values of $a_p$, for each fixed
combination of the remaining retained indices, equals exactly one.

% ------------------------------------------------------------------
%  TIKZ LATTICE
% ------------------------------------------------------------------
\begin{figure}[H]
\centering
\resizebox{0.95\linewidth}{!}{%
  % --- ADJUSTABLE PARAMETERS ---
\def\hsep{4.0}
\def\vsep{1.4}
\def\logicXshift{7.5}
% -----------------------------
\begin{tikzpicture}[
  every node/.style = {
    draw, rounded corners=3pt, inner sep=5pt,
    align=center, font=\small, thick
  },
  rootBox/.style   = {fill=blue!10, draw=blue!60!black},
  sumBox/.style    = {fill=orange!10, draw=orange!60!black},
  binBox/.style    = {fill=red!5, draw=red!80!black, dashed},
  arr/.style       = {->, >=stealth, thick},
  barr/.style      = {->, >=stealth, thick, dashed, red!80!black},
  lbl/.style       = {draw=none, fill=none, font=\footnotesize\itshape, inner sep=2pt},
  blbl/.style      = {draw=none, fill=none, font=\footnotesize\bfseries, inner sep=2pt, text=red!80!black},
  graylbl/.style   = {draw=none, fill=none, font=\small\bfseries, gray!50}
]

%% --- SUM AGGREGATION SECTION (Left) ---
\node[rootBox] (root) at (0,0) {$\phi_{cijm}$};
\node[sumBox] (Fcij)  at (-\hsep/2, \vsep)
  {$\widetilde{\phi}_{cij} \coloneqq \sum_{m \in \mathcal{M}} \phi_{cijm}$};
\node[sumBox] (Fijm)  at (\hsep/2, \vsep)
  {$\widetilde{\phi}_{ijm} \coloneqq \sum_{c \in C} \phi_{cijm}$};
\node[sumBox] (Fij)   at (0, 2*\vsep)
  {$\widetilde{\phi}_{ij} \coloneqq \sum_{c \in C} \sum_{m \in \mathcal{M}} \phi_{cijm}$};

\draw[arr] (root.north) -- (Fcij.south);
\draw[arr] (root.north) -- (Fijm.south);
\draw[arr] (Fcij.north) -- (Fij.south);
\draw[arr] (Fijm.north) -- (Fij.south);

% Dot notation
\node[sumBox] (dotj) at (-\hsep/2, -\vsep)
  {$\widetilde{\phi}_{c \cdot j m} \coloneqq \sum_{i \in V} \phi_{cijm}$};
\node[sumBox] (doti) at (\hsep/2, -\vsep)
  {$\widetilde{\phi}_{ci \cdot m} \coloneqq \sum_{j \in V} \phi_{cijm}$};

\draw[arr] (root.south) -- (dotj.north);
\draw[arr] (root.south) -- (doti.north);

%% --- BINARY LOGIC AGGREGATION SECTION (Right) ---
\begin{scope}[xshift=\logicXshift cm]
\node[rootBox] (Xroot) at (0, 0) {$X_{cijm}$};
\node[binBox] (XOR)  at (0, \vsep*1.36)
  {$\widecheck{X}_{cij}$ (OR aggregate) \\
   \scriptsize $\widecheck{X}_{cij} \leq \widetilde{X}_{cij}$ \\
   \scriptsize $\widetilde{X}_{cij} \leq |\mathcal{M}|\,\widecheck{X}_{cij}$};
\node[binBox] (XHAT) at (0, -\vsep*1.36)
  {$\widehat{X}_{cij}$ (At-most-one) \\
   \scriptsize $\widehat{X}_{cij} \leq \widetilde{X}_{cij}$ \\
   \scriptsize $\widetilde{X}_{cij} \leq \widehat{X}_{cij}$};

\draw[barr] (Xroot.north) -- (XOR.south)  node[blbl, midway, right] {Rule 3};
\draw[barr] (Xroot.south) -- (XHAT.north) node[blbl, midway, right] {Rule 4};
\end{scope}

%% --- DIVIDER AND LABELS ---
\draw[thick, gray!30, dashed]
  (\logicXshift/2, -1.5*\vsep) -- (\logicXshift/2, 2.8*\vsep);
\node[graylbl] at (0,            2.7*\vsep) {Sum Aggregation};
\node[graylbl] at (\logicXshift, 2.7*\vsep) {Binary Logic Aggregation};

\end{tikzpicture}
}
\caption{%
  Derivation lattice for a generic variable $\phi$ (left panel,
  continuous or integer) and its binary counterpart $X$ (right
  panel), separated by the dashed vertical rule.
  On the left, every arrow represents a sum aggregation
  (Rules~1--2): purely notational, no new variable introduced.
  On the right, solid arrows represent the same sum aggregation
  applied to binary variables, producing an integer count.
  Dashed arrows introduce a genuinely new solver variable (dashed
  border): the OR aggregate $\overline{X}$ (Rule~3), defined as a
  binary proxy for whether the count is positive, enforced by the
  big-$M$ pair shown in each node.
  The hat variant $\hat{X}$ (Rule~4) is the same variable as the
  corresponding OR aggregate with the additional
  constraint~\eqref{eq:exactly-one} in force.%
}
\label{fig:index-lattice}
\end{figure}

\end{convention}
